\begin{helper}[Third normal form in the last-index case]
Let \((\varphi,T)\) be a minimum bad pair with
\(2<\rho(T)=|E(T)|=n\).  Let the exact last edge have color
\(\gamma\), and suppose \(\gamma\) is present at the penultimate vertex.
There are an explicit replacement coloring \(\varphi_3\), a color
\(\beta\), and \(j\le n-2\) such that \((\varphi_3,T)\) is minimum bad,
the last edge still has color \(\gamma\), \(\gamma\) is still present at
the penultimate vertex, and \(\beta\) is missing at both \(y_j,y_n\).
\end{helper}
\begin{proof}
Unfolding \noderef{TashTreeMinimumBadPair} and
\noderef{TashTreeBadPair}, the full vertex set is not elementary and the
pair contains a \(k\)-triple; \noderef{KTriple} in particular supplies
properness of \(\varphi\).  The proper prefix through \(y_{n-1}\) is
elementary by \noderef{TashTreeProperPrefixElementary}.  By
\noderef{ColoredVertexSetElementary}, failure of elementariness on the
full set supplies two distinct vertices with a common missing color
\(\beta\).  The two cannot both lie in the elementary prefix, while
\noderef{TreeSequence} and \noderef{TreeSequencePrefix} show that the
only displayed vertex outside that prefix is \(y_n\).  Hence, in the
sense of \noderef{MissingColors}, \(\beta\) is missing at \(y_n\) and at
\(y_{j'}\) for some \(j'<n\).

If \(j'\le n-2\), use the original coloring.  Otherwise \(j'=n-1\).
Then \(\beta\) is unused on the colored edges of the prefix: if such an
edge had color \(\beta\), \noderef{TashkinovTree} would provide an
earlier missing-\(\beta\) witness.  That witness precedes the edge and is
therefore distinct from \(y_{n-1}\) by \noderef{TreeSequence}; together
with the missing-\(\beta\) status of \(y_{n-1}\), this contradicts
\noderef{ColoredVertexSetElementary} by the meaning of
\noderef{MissingColors}.  The last edge has endpoint \(y_n\) by
\noderef{TreeSequence}, so \noderef{MissingColors} says that it also
cannot have color \(\beta\).  Thus \(\beta\) is unused on the whole
displayed tree.

Apply \noderef{TashTreeFourUnusedColors} to the prefix through
\(y_{n-2}\).  By the meanings of \noderef{TreeFreeColors} and
\noderef{TreeUsedColors}, after excluding the at most two colors on the
two remaining exact edges there are at least two colors missing in the
prefix and unused on all displayed colored tree edges.  Choose one
\(\alpha\ne\beta\), with missing witness \(y_m\), \(m\le n-2\).
The two colors are unused through \(y_{n-1}\), so
\noderef{TashTreeControlledInterchange} applies at \(y_m,y_{n-1}\).
It returns an explicit component-swap coloring, makes \(\beta\) missing
at \(y_m\), and returns a minimum-bad certificate.

The last vertex \(y_n\), which misses \(\beta\), is a third
degree-at-most-one vertex by \noderef{MissingColorIncidenceEndpoint}:
the affected colors are distinct by choice, the coloring is proper by
\noderef{KTriple}, and \(\beta\) is missing at \(y_n\); apply the lemma
with \(\beta\) as its missing color.  If \(y_n\) lay on the controlled
component, \noderef{TwoColorComponentSpanningPath} would make it a third
end besides \(y_m,y_{n-1}\).  These three vertices are distinct by the
vertex-injectivity clause of \noderef{TreeSequence}, so this is
impossible.  The outside-component clause of
\noderef{TwoColorComponentSwapMissingBehavior} now preserves the whole
missing-color set at \(y_n\), in particular its missing-\(\beta\)
status.  The controlled-end conclusion of
\noderef{TashTreeControlledInterchange} makes \(\beta\) missing at
\(y_m\), and that theorem also supplies the asserted new minimum-bad
certificate.

Neither affected color equals \(\gamma\).  The color \(\alpha\) is
outside \noderef{TreeUsedColors} for the whole tree, while the displayed
last edge has color \(\gamma\), so \(\alpha\ne\gamma\).  Also, that last
edge is incident with \(y_n\) by \noderef{TreeSequence}; since it has
color \(\gamma\) while \(y_n\) misses \(\beta\),
\noderef{MissingColors} gives \(\beta\ne\gamma\).  The component swap
\noderef{TwoColorComponentSwap} changes only values equal to
\(\alpha\) or \(\beta\), and consequently fixes the exact last edge's
\(\gamma\)-value.  Finally,
\(\gamma\in\noderef{PresentColors}(\varphi,y_{n-1})\) supplies an exact
edge incident with \(y_{n-1}\) whose old value is \(\gamma\).  Whether
or not its edge-node lies in the swapped component, the cases of
\noderef{TwoColorComponentSwap} fix this value because
\(\gamma\notin\{\alpha,\beta\}\).  The same incident edge therefore
witnesses that \(\gamma\) remains present at \(y_{n-1}\).  Taking
\(j=m\) gives the replacement coloring and all required preserved
invariants.
\end{proof}
